\section{독자 계약과 전체 지도}
\label{sec:reading-contract}

\subsection{이 문서가 가정하는 것과 설명하는 것}

이 문서는 독자가 환(ring), 몫환(quotient ring), 자유 모듈(free module),
동치관계(equivalence relation)와 몫집합(quotient set),
가환도식(commutative diagram), 귀납법(induction)에 익숙하다고 가정한다.
수학 용어는 이처럼 문서의 첫 등장과 표제·도표·용어 사전에서 한국어 뒤에 영어 원어를
병기한다. 같은 문맥에서 반복할 때에는 가독성을 위해 한국어만 쓸 수 있다. 반대로
다음 지식은 가정하지 않는다.

\begin{itemize}
  \item HAETAE 또는 모듈 격자 서명(module-lattice signature)의 설계;
  \item Jasmin의 문법과 메모리 모델;
  \item EasyCrypt의 Hoare logic, pRHL, 전술 언어;
  \item rANS entropy coding의 구현 관례.
\end{itemize}

목표는 모든 소스 코드를 줄마다 설명하는 것이 아니다. 각 증명 의무에 대해
``대상(object)은 무엇인가'', ``어떤 사상(map)이 가환(commute)해야 하는가'',
``정의역(domain)을 제한하는 전제(premise)는 무엇인가'', ``현재 어느
화살표(arrow)까지 실제 절차에 대해 증명되었는가''를
답하는 것이 목표다. 프로젝트의 범위는 mode~2와 고정된 명세·구현 스냅샷에
한정된다 \cite{project-scope}.

\begin{translation}[세 층 번역]
저장소의 정리(theorem)는 다음 세 층을 오가는 명제(proposition)로 읽는다.
\begin{enumerate}
  \item \textbf{대수 층(algebraic layer)}: \(R_q\)-모듈(module) 원소,
        분해(decomposition), challenge, 논문 알고리즘;
  \item \textbf{표현 층(representation layer)}: 정준 바이트열(canonical byte
        string), 접두부 투영(prefix projection), 메모리 구간(memory region),
        고정폭 정수(fixed-width integer);
  \item \textbf{실행 층(execution layer)}: 실제 생성 절차, 분기, 루프,
        반환값과 trace 관측값(observation).
\end{enumerate}
한 층에서 성립하는 등식(equality)을 다른 층으로 올리거나 내리려면 반드시 명명된
representation/refinement 정리(theorem)가 필요하다. 같은 수학적 의도를 구현한다는
사실만으로 도식(diagram)이 자동으로 가환(commute)하지 않는다.
\end{translation}

\begin{figure}[ht]
\centering
\begin{tikzpicture}[node distance=8mm and 7mm]
  \node[layernode] (alg) {논문 대수 층\\(algebraic layer)\\\(R_q\)-모듈과 분포(distribution)};
  \node[layernode, right=of alg] (repr) {정준 표현 층\\(canonical representation)\\바이트열과 접두부};
  \node[layernode, right=of repr] (proc) {실제 실행 층\\(execution layer)\\Jasmin 상태기계};
  \node[partialnode, below=12mm of repr] (security) {최상위 기능·보안\\합성 목표};
  \draw[flowarrow] (alg) -- node[above,align=center]{encode/\\refine} (repr);
  \draw[flowarrow] (repr) -- node[above,align=center]{extract/\\simulate} (proc);
  \draw[openarrow] (proc) -- (security);
  \draw[openarrow] (alg) -- (security);
\end{tikzpicture}
\caption{전체 증명 구조. 수평 화살표의 일부만 닫혔고, 최상위 합성 화살표는
아직 부분적이다.}
\label{fig:three-layers}
\end{figure}

\subsection{상태어는 논리적 양화사(quantifier)를 포함한다}

운영상 source of truth는 \sourcepath{CLAIM_LEDGER.md}이며
\cite{claim-ledger}, 의존성은 \sourcepath{THEOREM_GRAPH.md}에 기록된다
\cite{theorem-graph}. 이 문서에서 상태어는 다음 의미로만 사용한다.

\begin{center}
\begin{tabularx}{\textwidth}{L{28mm} Y Y}
\toprule
표시 & 정확한 뜻 & 자동으로 따라오지 않는 것 \\
\midrule
\statusproved & 표시된 전제 아래 명명된 EasyCrypt 선언이 fresh-compile된다. &
부모 주장, 종료성(termination), losslessness,
분포 동일성(distributional equality) \\
\statuspartial & 자식 명제(proposition) 또는 제한된 방향은 증명되었지만 명명된 부모가 열려 있다. &
제한을 제거한 전역 명제(global statement) \\
\statusspecified & 목표를 나타내는 연산(operation)·술어(predicate)·인터페이스가
컴파일된다. &
그 목표의 증명 \\
\statusblocked & 구체적인 extraction, 의미론(semantics), 수치·확률 정리(theorem)가 없어
합성할 수 없다. &
난점이 해결 불가능하다는 주장 \\
\statusdeferred & 시간 제한 뒤 명시적인 논문 경계로 동결했다. &
거짓 또는 불필요하다는 주장 \\
\bottomrule
\end{tabularx}
\end{center}

특히 ``Hoare partial correctness가 증명되었다''는 문장은
\[
  \text{절차가 반환한다면 사후조건이 성립한다}
\]
는 뜻이다. 절차가 모든 입력에서 반환한다거나, 확률적 재시도가 거의 확실히
끝난다는 뜻이 아니다. 이 구분은 KeyGen과 Sign의 rejection loop 때문에
본질적이다.

\subsection{대수기하학은 어디까지 관련되는가}

\(R_q\)는 물론 아핀 스킴(affine scheme) \(\operatorname{Spec}R_q\)를
정의한다. 또한 메모리의 부분구간 제한(local restriction), 국소 등식(local
equality), frame 조건은 제한사상(restriction map)과 접합(gluing)을 떠올리게
한다. 그러나 현재 EasyCrypt 파일은 스킴(scheme), 다형체(variety),
스펙트럼(spectrum), 층(sheaf), 코호몰로지(cohomology)를 형식화하지 않는다.
증명 대상은 주로 유한 집합(finite set) 위의 배열, 고정폭 word,
상태 전이(state transition), 정수 나눗셈(integer division)이다.

\begin{analogy}
메모리 \(m\)을 주소집합(address set) 위의 단면(section)처럼 보고, 구간
\(U\)에서의 동일성을
\(m_1\restrict_U=m_2\restrict_U\)로 읽으면 region-local hash 정리(theorem)의 모양을
빠르게 이해할 수 있다. 하지만 이 비유에서 층 공리(sheaf axiom)나
스킴 구조(scheme structure)를 사용해 EasyCrypt 결론을 도출하지 않는다.
실제 근거는 점별 \identifier{load8}
등식과 절차의 read footprint다.
\end{analogy}

대수기하학자가 특히 주의할 점은 hash 함수다. SHAKE 상태전이는 일반적으로
\(R_q\)-선형사상(linear map)도 정칙사상(regular morphism)도 아니다. 이
문서에서 hash는 유한
바이트열을 읽는 결정론적 상태기계로 취급한다. 따라서 hash 입력의 동일성에서
출력 동일성(output equality)을 얻는 이유는 대수적 함수성(algebraic
functionality)보다 실제 두 프로그램의 관계적 실행 정리(relational execution
theorem)에 있다.

\subsection{먼저 기억할 여섯 문장}

\begin{enumerate}
  \item 최상위 \identifier{FC_KeyGen}, \identifier{FC_Sign},
        \identifier{FC_Verify}는 아직 목표 인터페이스이지 완료된 정리(theorem)가 아니다.
  \item \(\prefix_{992}(sk)=vk\)와 raw/internal \(\mu\)-접두부 일치는 실제
        생성 절차까지 증명되었지만, 전체 API 보안으로 자동 승격되지 않는다.
  \item HBZ/rANS의 순수 역함수 구조, 실제 encoder/decoder의 개별 trace
        정제(refinement), actual copy를 포함한 core harness 역함수(inverse
        function), production full-HBZ wrapper 화살표까지 성공 조건부
        부분정확성(success-conditioned partial correctness)으로 닫혔다.
  \item canonical all-zero HBZ가 만드는 all-6 stream에서는 반환하는 실제
        encoder의 실패가 배제되고 production zero round trip도 증명되었지만,
        고정 입력의 종료성(termination), losslessness와 \identifier{phoare}는
        열려 있다.
  \item actual KeyGen matrix/finalize snapshot에서는 KG-2-like 분해와
        snapshot-only mod-\(2q\) 합동식이 증명되었지만,
        마지막 \identifier{output_row} NTT-word rewrite 뒤의 full-NTT convolution과
        보안 모형 표현 어댑터가 없어 논문 KeyGen 식으로 승격되지 않는다.
        따라서 KeyGen은 \identifier{STOP-KG-NTT}로 동결되었고, 다음 actual
        accepted Sign core는 아직 \statusspecified 상태다.
  \item 최종 암호학적 가정과 구현 증명 의무를 분리한다. MLWE/BST-MSIS/ROM
        인터페이스 외의 구현 정확성은 가정으로 숨기지 않는다.
\end{enumerate}

이 여섯 문장은 이후 모든 장의 상태 판정 규칙이다.
